%%%%%%%%%%%%%%%%%%%%%%%%%%%%%%%%%%%%%%%%%%%%%%%%%%%%%%%%%%%%%%%%%%%%%%
% Lớp 10 | Véctơ (chưa xét tọa độ) | Tích của một số với véctơ
% Dạng 0H5?3-3: Tìm điểm nhờ đẳng thức véctơ
% Số câu: 2  (sinh tự động bởi tach_cau.py; từ nay thầy sửa trực tiếp file này)
%%%%%%%%%%%%%%%%%%%%%%%%%%%%%%%%%%%%%%%%%%%%%%%%%%%%%%%%%%%%%%%%%%%%%%

%%%%%%%%%%%%%%%% ĐÚNG/SAI %%%%%%%%%%%%%%%%

% nguồn: MaHoaNH_DungSai_Lop10
\begin{ex}%[0H5N3-3]%[ma:b9823f41]
	Cho hai điểm phân biệt $A$  và $B$. Hãy xét tính đúng sai của các mệnh đề: Điều kiện để điểm $I$ là trung điểm của đoạn thẳng $AB$ là
	 \choiceTF
		{$IA=IB$} 
		{$\overrightarrow{IA}=\overrightarrow{IB}$} 
		{\True $\overrightarrow{IA}=-\overrightarrow{IB}$} 
		{$\overrightarrow{AI}=\overrightarrow{BI}$} 
	\loigiai{ 
		\begin{itemchoice}
			\itemch $IA=IB$ khi và chỉ khi $I$ thuộc đường trung  trực của đoạn thẳng $AB$.
			\itemch $\overrightarrow{IA}=\overrightarrow{IB}$  khi và chỉ khi $\overrightarrow{BA}=\vec{0}$. Khi đó $A$ trùng với $B$.
			\itemch $\overrightarrow{IA}=-\overrightarrow{IB}$ khi và chỉ khi $I$ là trung điểm của đoạn thẳng  $AB$.
			\itemch $\overrightarrow{AI}=\overrightarrow{BI}$ khi và chỉ khi $\overrightarrow{AB}=\vec{0}$. Khi đó $A$ trùng  với $B$.
		\end{itemchoice}
	} 
\end{ex}

% nguồn: MaHoaNH_DungSai_Lop10
\begin{ex}%[0H5H3-3]%[ma:a81ed701]
	Cho tam giác $ABC$ có $G$ là trọng tâm và $I$ là trung điểm của đoạn thẳng $BC$. Xét tính đúng sai của các  mệnh đề sau:
	\choiceTF 
		{$\overrightarrow{GA}=2\overrightarrow{GI}$} 
		 {$\overrightarrow{IG}=-\dfrac{1}{3}\overrightarrow{IA}$}
		 {\True $\overrightarrow{GB}+\overrightarrow{GC}=2\overrightarrow{GI}$}
		 {$\overrightarrow{GB}+\overrightarrow{GC}=\overrightarrow{GA}$}
	 \loigiai{
		\begin{itemchoice}
			\itemch Theo tính chất trọng tâm, ta  có $\overrightarrow{GA}=-2\overrightarrow{GI}$.
			\itemch Theo tính chất trọng  tâm, ta có $\overrightarrow{IG}=\dfrac{1}{3}\overrightarrow{IA}$
			\itemch Do $I$ là trung điểm của  $BC$ nên $\overrightarrow{GB}+\overrightarrow{GC}=2\overrightarrow{GI}$
			\itemch Theo tính chất trọng tâm, ta có $\overrightarrow{GB}+\overrightarrow{GC}+\overrightarrow{GA}=\vec{0}$, suy  ra $\overrightarrow{GB}+\overrightarrow{GC}=-\overrightarrow{GA}$.
		\end{itemchoice}
	} 
\end{ex}
